% =====================================================================
% EE 414 - Introduction to Robotics
% Shared beamer preamble for the weekly lecture decks.
%
% A deck is:
%     \documentclass[aspectratio=169]{beamer}
%     \usetheme{Madrid}\usecolortheme{default}
%     % =====================================================================
% EE 414 - Introduction to Robotics
% Shared beamer preamble for the weekly lecture decks.
%
% A deck is:
%     \documentclass[aspectratio=169]{beamer}
%     \usetheme{Madrid}\usecolortheme{default}
%     % =====================================================================
% EE 414 - Introduction to Robotics
% Shared beamer preamble for the weekly lecture decks.
%
% A deck is:
%     \documentclass[aspectratio=169]{beamer}
%     \usetheme{Madrid}\usecolortheme{default}
%     % =====================================================================
% EE 414 - Introduction to Robotics
% Shared beamer preamble for the weekly lecture decks.
%
% A deck is:
%     \documentclass[aspectratio=169]{beamer}
%     \usetheme{Madrid}\usecolortheme{default}
%     \input{../../shared/ee414-beamer-preamble.tex}
%     \title...\author...\date...
%     \begin{document} ... \end{document}
%
% Build any deck with TWO pdflatex passes (three if it uses \fullslide,
% whose TikZ page anchors come from the .aux).
%
% Adapted from the SE 100 preamble: same palette, same terminal-window code
% rule, same "under the lid" device renamed \begin{underhood} for this course.
% =====================================================================

\usepackage[T1]{fontenc}
\usepackage[utf8]{inputenc}
\usepackage{graphicx}
\usepackage{amsmath}
\usepackage{amssymb}
\usepackage{tikz}
\usetikzlibrary{arrows.meta,shapes.geometric,positioning,calc,fit,backgrounds}
\usepackage{booktabs}
\usepackage{array}
\usepackage{hyperref}

\graphicspath{{./}{./figures/}}

% ---------------------------------------------------------------
% House colour palette - shared with the other Alfaisal course decks.
\definecolor{primary}{RGB}{10,37,64}
\definecolor{secondary}{RGB}{99,91,255}
\definecolor{accent}{RGB}{0,212,255}
\definecolor{success}{RGB}{50,213,131}
\definecolor{warning}{RGB}{255,193,7}
\definecolor{danger}{RGB}{220,53,69}
\definecolor{softgrey}{RGB}{240,242,245}

\setbeamercolor{structure}{fg=secondary}
\setbeamercolor{palette primary}{bg=primary,fg=white}
\setbeamercolor{palette secondary}{bg=secondary,fg=white}
\setbeamercolor{palette tertiary}{bg=accent,fg=white}
\setbeamertemplate{navigation symbols}{}
\setbeamerfont{block title}{size=\normalsize}

\AtBeginSection[]{
  \begin{frame}[plain]
    \vfill
    \centering
    \begin{beamercolorbox}[sep=10pt,center,rounded=true,shadow=true]{palette primary}
      \usebeamerfont{title}\insertsectionhead
    \end{beamercolorbox}
    \vfill
  \end{frame}
}

% ---------------------------------------------------------------
% TikZ styles.
%   sense / plan / act  - the three stages of the autonomy loop, colour-coded
%                         once here and used identically in every deck, so a
%                         student reads the colour before reading the label.
%   node2 / topicbox    - ROS 2 computation-graph drawings.
\tikzset{
  sense/.style={rectangle, rounded corners=4pt, minimum width=2.5cm,
                minimum height=0.8cm, text centered, draw=accent!80!black,
                fill=accent!15, font=\footnotesize},
  plan/.style={rectangle, rounded corners=4pt, minimum width=2.5cm,
               minimum height=0.8cm, text centered, draw=secondary,
               fill=secondary!12, font=\footnotesize},
  act/.style={rectangle, rounded corners=4pt, minimum width=2.5cm,
              minimum height=0.8cm, text centered, draw=success!70!black,
              fill=success!18, font=\footnotesize},
  world/.style={ellipse, minimum width=2.4cm, minimum height=0.9cm,
                text centered, draw=primary, fill=primary!10, font=\footnotesize},
  rosnode/.style={ellipse, minimum width=2.4cm, minimum height=0.85cm,
                  text centered, draw=secondary, fill=secondary!12,
                  font=\footnotesize\ttfamily},
  topicbox/.style={rectangle, minimum width=2.0cm, minimum height=0.55cm,
                   text centered, draw=primary!60, fill=softgrey,
                   font=\scriptsize\ttfamily},
  hw/.style={rectangle, minimum width=2.2cm, minimum height=0.8cm,
             text centered, draw=warning!80!black, fill=warning!20,
             font=\footnotesize},
  flowarrow/.style={-{Stealth[length=2.5mm]}, thick, draw=primary},
  dataflow/.style={-{Stealth[length=2mm]}, semithick, draw=secondary}
}

% ---------------------------------------------------------------
% Full-bleed infographic slide. The title bar is part of the image,
% so the frame carries no beamer title and no furniture.
\newcommand{\fullslide}[1]{%
  \begin{frame}[plain]
    \begin{tikzpicture}[remember picture,overlay]
      \node at (current page.center)
        {\includegraphics[width=\paperwidth,height=\paperheight]{#1}};
    \end{tikzpicture}
  \end{frame}%
}

% ---------------------------------------------------------------
% RULE: every code box in this course wears the macOS terminal window bar.
% The three traffic lights are what tell a student, at a glance, that they are
% looking at something that RUNS ON A MACHINE rather than at prose about it.
% There is no bare-box variant: the title keys are applied AFTER the caller's
% options, so passing [notitle] has no effect. Do not add an exception.
%
%   \begin{pycode} ... \end{pycode}   a Python (rclpy) source file
%   \begin{shell}  ... \end{shell}    a terminal session or program output
%   \begin{spec}   ... \end{spec}     NOT runnable: .msg definitions, URDF
%                                     fragments, pseudocode, mathematics
%
% The `...s' variants (pycodes, shells) are the small ones, for a crowded slide.
%
% STANDING CONSTRAINT -- read this before writing a .msg or a shell comment.
% Inside ANY of these boxes, a `#' MUST NOT be the first non-space character of a
% line. Beamer's [fragile] scanner re-reads the frame body and treats a leading
% `#' as a macro parameter, which fails with
%       ! Illegal parameter number in definition of \next.
% Leading whitespace does not help -- it is stripped before the scan.
% A `#' anywhere else on the line is fine, which is why inline comments work:
%       self.create_timer(0.1, self.tick)   # 10 Hz, forever      <- OK
%       # geometry_msgs/msg/Twist                                 <- BREAKS THE BUILD
% For a header comment, put the file or type name in the prose ABOVE the box.
\usepackage{tcolorbox}
\tcbuselibrary{listings,skins}
\usepackage{listings}

\definecolor{termbg}{RGB}{30,30,32}
\definecolor{termbar}{RGB}{58,58,60}
\definecolor{termedge}{RGB}{88,88,92}
\definecolor{termfg}{RGB}{219,219,219}
\definecolor{termcom}{RGB}{106,153,85}
\definecolor{termkw}{RGB}{86,156,214}
\definecolor{termstr}{RGB}{206,145,120}
\definecolor{termfun}{RGB}{220,220,170}
\definecolor{termerr}{RGB}{241,76,76}
\definecolor{dotr}{RGB}{255,95,86}
\definecolor{doty}{RGB}{255,189,46}
\definecolor{dotg}{RGB}{39,201,63}

\newcommand{\termdots}{%
  \begin{tikzpicture}[baseline=-0.65ex]
    \fill[dotr] (0,0)     circle (2.7pt);
    \fill[doty] (0.34,0)  circle (2.7pt);
    \fill[dotg] (0.68,0)  circle (2.7pt);
  \end{tikzpicture}}

% Python source: VS Code "Dark+" palette, so what a student sees on the slide
% is what their editor shows them ten minutes later.
\lstdefinestyle{darkpy}{
  language=Python,
  basicstyle=\ttfamily\footnotesize\color{termfg},
  keywordstyle=\color{termkw},
  stringstyle=\color{termstr},
  commentstyle=\color{termcom},
  emph={print,len,range,float,int,str,super,min,max,abs,round},
  emphstyle=\color{termfun},
  emph={[2]rclpy,Node,create_publisher,create_subscription,create_timer,spin,
        Twist,LaserScan,Odometry,String},
  emphstyle={[2]\color{termfun}},
  showstringspaces=false, columns=fullflexible, keepspaces=true,
}
% Terminal output: same lexer, keywords left plain, errors picked out in red.
\lstdefinestyle{darkterm}{
  language=Python,
  basicstyle=\ttfamily\footnotesize\color{termfg},
  keywordstyle=\color{termfg},
  stringstyle=\color{termstr},
  commentstyle=\color{termcom},
  emph={ros2,colcon,rviz2,gazebo,source,export},
  emphstyle=\color{termfun},
  emph={[2]ERROR,error,FATAL,Failed,failed,Traceback},
  emphstyle={[2]\color{termerr}\bfseries},
  emph={[3]INFO,OK,PASS},
  emphstyle={[3]\color{termcom}},
  showstringspaces=false, columns=fullflexible, keepspaces=true,
}

\tcbset{termbox/.style={
  enhanced, listing only,
  colback=termbg, colframe=termedge, boxrule=0.5pt,
  arc=9pt, outer arc=9pt,
  left=10pt, right=10pt, top=6pt, bottom=6pt,
  title={\termdots}, colbacktitle=termbar, coltitle=termfg,
  toptitle=3.5pt, bottomtitle=3.5pt, lefttitle=10pt, titlerule=0pt,
}}
\lstdefinestyle{darkpys}{style=darkpy,     basicstyle=\ttfamily\scriptsize\color{termfg}}
\lstdefinestyle{darkterms}{style=darkterm, basicstyle=\ttfamily\scriptsize\color{termfg}}
\tcbset{tightpad/.style={top=4pt, bottom=4pt, toptitle=3pt, bottomtitle=3pt}}

\newtcblisting{pycode}[1][]{termbox, listing style=darkpy,    #1, title={\termdots}, colbacktitle=termbar, coltitle=termfg}
\newtcblisting{shell}[1][]{termbox,  listing style=darkterm,  #1, title={\termdots}, colbacktitle=termbar, coltitle=termfg}
\newtcblisting{pycodes}[1][]{termbox, listing style=darkpys,  tightpad, #1, title={\termdots}, colbacktitle=termbar, coltitle=termfg}
\newtcblisting{shells}[1][]{termbox,  listing style=darkterms, tightpad, #1, title={\termdots}, colbacktitle=termbar, coltitle=termfg}

% The light box: for things that are NOT runnable -- .msg and .srv interface
% definitions, URDF fragments, pseudocode. The contrast with the terminal box is
% the point: no window bar means nothing here runs.
%
% NOTE: the style MUST be a named \lstdefinestyle used via `listing style='.
% Passing `listing options={...}' inline instead re-parses the body and breaks on
% any `#' in the listing -- which every .msg comment starts with.
\lstdefinestyle{specstyle}{
  basicstyle=\ttfamily\footnotesize\color{primary!85!black},
  commentstyle=\color{primary!55},
  columns=fullflexible, keepspaces=true, showstringspaces=false,
}
\newtcblisting{spec}[1][]{
  enhanced, listing only, listing style=specstyle,
  colback=softgrey, colframe=primary!35, boxrule=0.6pt, arc=3pt,
  left=8pt, right=8pt, top=5pt, bottom=5pt, #1}

% ---------------------------------------------------------------
% "Under the hood" - the machine-level aside.
%
% STANDING DEVICE. Wherever a slide states a rule about how a robot or ROS 2
% behaves, the reason belongs in one of these, immediately after it. A student
% who can answer "why does it behave that way?" with something other than
% "because ROS does" is not memorising.
%
% FOUR LINES MAXIMUM. Deliberately out of scope in this course: the DDS wire
% protocol, real-time kernel scheduling, and the internals of the executor.
% Naming the exclusion here is what stops it creeping in.
\newenvironment{underhood}[1][Under the hood]{%
  \begin{tcolorbox}[enhanced, colback=primary!6, colframe=primary!45,
                    boxrule=0.6pt, arc=3pt,
                    left=8pt, right=8pt, top=5pt, bottom=5pt,
                    title={#1}, colbacktitle=primary!85!black, coltitle=white,
                    fonttitle=\footnotesize\bfseries,
                    toptitle=2pt, bottomtitle=2pt, titlerule=0pt]%
}{\end{tcolorbox}}

% "In the field" - the industrial-reality aside. Same shape, different colour:
% where \underhood explains a mechanism, this one says what it costs in a real
% deployment. Use sparingly, and never as decoration.
\newenvironment{inthefield}[1][In the field]{%
  \begin{tcolorbox}[enhanced, colback=success!7, colframe=success!50!black,
                    boxrule=0.6pt, arc=3pt,
                    left=8pt, right=8pt, top=5pt, bottom=5pt,
                    title={#1}, colbacktitle=success!45!black, coltitle=white,
                    fonttitle=\footnotesize\bfseries,
                    toptitle=2pt, bottomtitle=2pt, titlerule=0pt]%
}{\end{tcolorbox}}

% ---------------------------------------------------------------
% \imageslide{title}{figure file}{the one sentence to say about it}
%
% A PDF shows video as a still, so the deck uses images. The clip, if there is
% one, is opened from a browser -- its URL lives in the week README, never on a
% slide, because a dead link projected in front of a class is worse than no link.
%
% RULE: an image slot carries ONE sentence, and that sentence names what the
% students are looking at in the course's own vocabulary. A picture with no
% naming sentence is decoration.
%
% Generated figures are made with ../shared/gen_figure.py and each keeps its
% prompt beside it as <name>.prompt.txt, so any figure in this course can be
% regenerated or restyled consistently.
\newcommand{\imageslide}[3]{%
  \begin{frame}{#1}
    \vspace{-0.2em}
    \begin{center}
      \includegraphics[width=0.93\textwidth,height=0.63\textheight,keepaspectratio]{#2}
    \end{center}
    \vspace{0.3em}
    \centering{\small #3}
  \end{frame}%
}

% ---------------------------------------------------------------
% THE PINNED STACK. This macro and setup/README.md are the ONLY two places in
% the course where the distribution is named. Every deck, exercise and lab
% sheet writes \rosver, never the literal name. A distribution upgrade is then
% two edits, not a search across fifteen weeks of LaTeX that will miss three.
\newcommand{\rosver}{ROS\,2 Jazzy}
\newcommand{\ubuntuver}{Ubuntu 24.04 LTS}
\newcommand{\gazebover}{Gazebo Harmonic}
\newcommand{\pyver}{3.12}

%     \title...\author...\date...
%     \begin{document} ... \end{document}
%
% Build any deck with TWO pdflatex passes (three if it uses \fullslide,
% whose TikZ page anchors come from the .aux).
%
% Adapted from the SE 100 preamble: same palette, same terminal-window code
% rule, same "under the lid" device renamed \begin{underhood} for this course.
% =====================================================================

\usepackage[T1]{fontenc}
\usepackage[utf8]{inputenc}
\usepackage{graphicx}
\usepackage{amsmath}
\usepackage{amssymb}
\usepackage{tikz}
\usetikzlibrary{arrows.meta,shapes.geometric,positioning,calc,fit,backgrounds}
\usepackage{booktabs}
\usepackage{array}
\usepackage{hyperref}

\graphicspath{{./}{./figures/}}

% ---------------------------------------------------------------
% House colour palette - shared with the other Alfaisal course decks.
\definecolor{primary}{RGB}{10,37,64}
\definecolor{secondary}{RGB}{99,91,255}
\definecolor{accent}{RGB}{0,212,255}
\definecolor{success}{RGB}{50,213,131}
\definecolor{warning}{RGB}{255,193,7}
\definecolor{danger}{RGB}{220,53,69}
\definecolor{softgrey}{RGB}{240,242,245}

\setbeamercolor{structure}{fg=secondary}
\setbeamercolor{palette primary}{bg=primary,fg=white}
\setbeamercolor{palette secondary}{bg=secondary,fg=white}
\setbeamercolor{palette tertiary}{bg=accent,fg=white}
\setbeamertemplate{navigation symbols}{}
\setbeamerfont{block title}{size=\normalsize}

\AtBeginSection[]{
  \begin{frame}[plain]
    \vfill
    \centering
    \begin{beamercolorbox}[sep=10pt,center,rounded=true,shadow=true]{palette primary}
      \usebeamerfont{title}\insertsectionhead
    \end{beamercolorbox}
    \vfill
  \end{frame}
}

% ---------------------------------------------------------------
% TikZ styles.
%   sense / plan / act  - the three stages of the autonomy loop, colour-coded
%                         once here and used identically in every deck, so a
%                         student reads the colour before reading the label.
%   node2 / topicbox    - ROS 2 computation-graph drawings.
\tikzset{
  sense/.style={rectangle, rounded corners=4pt, minimum width=2.5cm,
                minimum height=0.8cm, text centered, draw=accent!80!black,
                fill=accent!15, font=\footnotesize},
  plan/.style={rectangle, rounded corners=4pt, minimum width=2.5cm,
               minimum height=0.8cm, text centered, draw=secondary,
               fill=secondary!12, font=\footnotesize},
  act/.style={rectangle, rounded corners=4pt, minimum width=2.5cm,
              minimum height=0.8cm, text centered, draw=success!70!black,
              fill=success!18, font=\footnotesize},
  world/.style={ellipse, minimum width=2.4cm, minimum height=0.9cm,
                text centered, draw=primary, fill=primary!10, font=\footnotesize},
  rosnode/.style={ellipse, minimum width=2.4cm, minimum height=0.85cm,
                  text centered, draw=secondary, fill=secondary!12,
                  font=\footnotesize\ttfamily},
  topicbox/.style={rectangle, minimum width=2.0cm, minimum height=0.55cm,
                   text centered, draw=primary!60, fill=softgrey,
                   font=\scriptsize\ttfamily},
  hw/.style={rectangle, minimum width=2.2cm, minimum height=0.8cm,
             text centered, draw=warning!80!black, fill=warning!20,
             font=\footnotesize},
  flowarrow/.style={-{Stealth[length=2.5mm]}, thick, draw=primary},
  dataflow/.style={-{Stealth[length=2mm]}, semithick, draw=secondary}
}

% ---------------------------------------------------------------
% Full-bleed infographic slide. The title bar is part of the image,
% so the frame carries no beamer title and no furniture.
\newcommand{\fullslide}[1]{%
  \begin{frame}[plain]
    \begin{tikzpicture}[remember picture,overlay]
      \node at (current page.center)
        {\includegraphics[width=\paperwidth,height=\paperheight]{#1}};
    \end{tikzpicture}
  \end{frame}%
}

% ---------------------------------------------------------------
% RULE: every code box in this course wears the macOS terminal window bar.
% The three traffic lights are what tell a student, at a glance, that they are
% looking at something that RUNS ON A MACHINE rather than at prose about it.
% There is no bare-box variant: the title keys are applied AFTER the caller's
% options, so passing [notitle] has no effect. Do not add an exception.
%
%   \begin{pycode} ... \end{pycode}   a Python (rclpy) source file
%   \begin{shell}  ... \end{shell}    a terminal session or program output
%   \begin{spec}   ... \end{spec}     NOT runnable: .msg definitions, URDF
%                                     fragments, pseudocode, mathematics
%
% The `...s' variants (pycodes, shells) are the small ones, for a crowded slide.
%
% STANDING CONSTRAINT -- read this before writing a .msg or a shell comment.
% Inside ANY of these boxes, a `#' MUST NOT be the first non-space character of a
% line. Beamer's [fragile] scanner re-reads the frame body and treats a leading
% `#' as a macro parameter, which fails with
%       ! Illegal parameter number in definition of \next.
% Leading whitespace does not help -- it is stripped before the scan.
% A `#' anywhere else on the line is fine, which is why inline comments work:
%       self.create_timer(0.1, self.tick)   # 10 Hz, forever      <- OK
%       # geometry_msgs/msg/Twist                                 <- BREAKS THE BUILD
% For a header comment, put the file or type name in the prose ABOVE the box.
\usepackage{tcolorbox}
\tcbuselibrary{listings,skins}
\usepackage{listings}

\definecolor{termbg}{RGB}{30,30,32}
\definecolor{termbar}{RGB}{58,58,60}
\definecolor{termedge}{RGB}{88,88,92}
\definecolor{termfg}{RGB}{219,219,219}
\definecolor{termcom}{RGB}{106,153,85}
\definecolor{termkw}{RGB}{86,156,214}
\definecolor{termstr}{RGB}{206,145,120}
\definecolor{termfun}{RGB}{220,220,170}
\definecolor{termerr}{RGB}{241,76,76}
\definecolor{dotr}{RGB}{255,95,86}
\definecolor{doty}{RGB}{255,189,46}
\definecolor{dotg}{RGB}{39,201,63}

\newcommand{\termdots}{%
  \begin{tikzpicture}[baseline=-0.65ex]
    \fill[dotr] (0,0)     circle (2.7pt);
    \fill[doty] (0.34,0)  circle (2.7pt);
    \fill[dotg] (0.68,0)  circle (2.7pt);
  \end{tikzpicture}}

% Python source: VS Code "Dark+" palette, so what a student sees on the slide
% is what their editor shows them ten minutes later.
\lstdefinestyle{darkpy}{
  language=Python,
  basicstyle=\ttfamily\footnotesize\color{termfg},
  keywordstyle=\color{termkw},
  stringstyle=\color{termstr},
  commentstyle=\color{termcom},
  emph={print,len,range,float,int,str,super,min,max,abs,round},
  emphstyle=\color{termfun},
  emph={[2]rclpy,Node,create_publisher,create_subscription,create_timer,spin,
        Twist,LaserScan,Odometry,String},
  emphstyle={[2]\color{termfun}},
  showstringspaces=false, columns=fullflexible, keepspaces=true,
}
% Terminal output: same lexer, keywords left plain, errors picked out in red.
\lstdefinestyle{darkterm}{
  language=Python,
  basicstyle=\ttfamily\footnotesize\color{termfg},
  keywordstyle=\color{termfg},
  stringstyle=\color{termstr},
  commentstyle=\color{termcom},
  emph={ros2,colcon,rviz2,gazebo,source,export},
  emphstyle=\color{termfun},
  emph={[2]ERROR,error,FATAL,Failed,failed,Traceback},
  emphstyle={[2]\color{termerr}\bfseries},
  emph={[3]INFO,OK,PASS},
  emphstyle={[3]\color{termcom}},
  showstringspaces=false, columns=fullflexible, keepspaces=true,
}

\tcbset{termbox/.style={
  enhanced, listing only,
  colback=termbg, colframe=termedge, boxrule=0.5pt,
  arc=9pt, outer arc=9pt,
  left=10pt, right=10pt, top=6pt, bottom=6pt,
  title={\termdots}, colbacktitle=termbar, coltitle=termfg,
  toptitle=3.5pt, bottomtitle=3.5pt, lefttitle=10pt, titlerule=0pt,
}}
\lstdefinestyle{darkpys}{style=darkpy,     basicstyle=\ttfamily\scriptsize\color{termfg}}
\lstdefinestyle{darkterms}{style=darkterm, basicstyle=\ttfamily\scriptsize\color{termfg}}
\tcbset{tightpad/.style={top=4pt, bottom=4pt, toptitle=3pt, bottomtitle=3pt}}

\newtcblisting{pycode}[1][]{termbox, listing style=darkpy,    #1, title={\termdots}, colbacktitle=termbar, coltitle=termfg}
\newtcblisting{shell}[1][]{termbox,  listing style=darkterm,  #1, title={\termdots}, colbacktitle=termbar, coltitle=termfg}
\newtcblisting{pycodes}[1][]{termbox, listing style=darkpys,  tightpad, #1, title={\termdots}, colbacktitle=termbar, coltitle=termfg}
\newtcblisting{shells}[1][]{termbox,  listing style=darkterms, tightpad, #1, title={\termdots}, colbacktitle=termbar, coltitle=termfg}

% The light box: for things that are NOT runnable -- .msg and .srv interface
% definitions, URDF fragments, pseudocode. The contrast with the terminal box is
% the point: no window bar means nothing here runs.
%
% NOTE: the style MUST be a named \lstdefinestyle used via `listing style='.
% Passing `listing options={...}' inline instead re-parses the body and breaks on
% any `#' in the listing -- which every .msg comment starts with.
\lstdefinestyle{specstyle}{
  basicstyle=\ttfamily\footnotesize\color{primary!85!black},
  commentstyle=\color{primary!55},
  columns=fullflexible, keepspaces=true, showstringspaces=false,
}
\newtcblisting{spec}[1][]{
  enhanced, listing only, listing style=specstyle,
  colback=softgrey, colframe=primary!35, boxrule=0.6pt, arc=3pt,
  left=8pt, right=8pt, top=5pt, bottom=5pt, #1}

% ---------------------------------------------------------------
% "Under the hood" - the machine-level aside.
%
% STANDING DEVICE. Wherever a slide states a rule about how a robot or ROS 2
% behaves, the reason belongs in one of these, immediately after it. A student
% who can answer "why does it behave that way?" with something other than
% "because ROS does" is not memorising.
%
% FOUR LINES MAXIMUM. Deliberately out of scope in this course: the DDS wire
% protocol, real-time kernel scheduling, and the internals of the executor.
% Naming the exclusion here is what stops it creeping in.
\newenvironment{underhood}[1][Under the hood]{%
  \begin{tcolorbox}[enhanced, colback=primary!6, colframe=primary!45,
                    boxrule=0.6pt, arc=3pt,
                    left=8pt, right=8pt, top=5pt, bottom=5pt,
                    title={#1}, colbacktitle=primary!85!black, coltitle=white,
                    fonttitle=\footnotesize\bfseries,
                    toptitle=2pt, bottomtitle=2pt, titlerule=0pt]%
}{\end{tcolorbox}}

% "In the field" - the industrial-reality aside. Same shape, different colour:
% where \underhood explains a mechanism, this one says what it costs in a real
% deployment. Use sparingly, and never as decoration.
\newenvironment{inthefield}[1][In the field]{%
  \begin{tcolorbox}[enhanced, colback=success!7, colframe=success!50!black,
                    boxrule=0.6pt, arc=3pt,
                    left=8pt, right=8pt, top=5pt, bottom=5pt,
                    title={#1}, colbacktitle=success!45!black, coltitle=white,
                    fonttitle=\footnotesize\bfseries,
                    toptitle=2pt, bottomtitle=2pt, titlerule=0pt]%
}{\end{tcolorbox}}

% ---------------------------------------------------------------
% \imageslide{title}{figure file}{the one sentence to say about it}
%
% A PDF shows video as a still, so the deck uses images. The clip, if there is
% one, is opened from a browser -- its URL lives in the week README, never on a
% slide, because a dead link projected in front of a class is worse than no link.
%
% RULE: an image slot carries ONE sentence, and that sentence names what the
% students are looking at in the course's own vocabulary. A picture with no
% naming sentence is decoration.
%
% Generated figures are made with ../shared/gen_figure.py and each keeps its
% prompt beside it as <name>.prompt.txt, so any figure in this course can be
% regenerated or restyled consistently.
\newcommand{\imageslide}[3]{%
  \begin{frame}{#1}
    \vspace{-0.2em}
    \begin{center}
      \includegraphics[width=0.93\textwidth,height=0.63\textheight,keepaspectratio]{#2}
    \end{center}
    \vspace{0.3em}
    \centering{\small #3}
  \end{frame}%
}

% ---------------------------------------------------------------
% THE PINNED STACK. This macro and setup/README.md are the ONLY two places in
% the course where the distribution is named. Every deck, exercise and lab
% sheet writes \rosver, never the literal name. A distribution upgrade is then
% two edits, not a search across fifteen weeks of LaTeX that will miss three.
\newcommand{\rosver}{ROS\,2 Jazzy}
\newcommand{\ubuntuver}{Ubuntu 24.04 LTS}
\newcommand{\gazebover}{Gazebo Harmonic}
\newcommand{\pyver}{3.12}

%     \title...\author...\date...
%     \begin{document} ... \end{document}
%
% Build any deck with TWO pdflatex passes (three if it uses \fullslide,
% whose TikZ page anchors come from the .aux).
%
% Adapted from the SE 100 preamble: same palette, same terminal-window code
% rule, same "under the lid" device renamed \begin{underhood} for this course.
% =====================================================================

\usepackage[T1]{fontenc}
\usepackage[utf8]{inputenc}
\usepackage{graphicx}
\usepackage{amsmath}
\usepackage{amssymb}
\usepackage{tikz}
\usetikzlibrary{arrows.meta,shapes.geometric,positioning,calc,fit,backgrounds}
\usepackage{booktabs}
\usepackage{array}
\usepackage{hyperref}

\graphicspath{{./}{./figures/}}

% ---------------------------------------------------------------
% House colour palette - shared with the other Alfaisal course decks.
\definecolor{primary}{RGB}{10,37,64}
\definecolor{secondary}{RGB}{99,91,255}
\definecolor{accent}{RGB}{0,212,255}
\definecolor{success}{RGB}{50,213,131}
\definecolor{warning}{RGB}{255,193,7}
\definecolor{danger}{RGB}{220,53,69}
\definecolor{softgrey}{RGB}{240,242,245}

\setbeamercolor{structure}{fg=secondary}
\setbeamercolor{palette primary}{bg=primary,fg=white}
\setbeamercolor{palette secondary}{bg=secondary,fg=white}
\setbeamercolor{palette tertiary}{bg=accent,fg=white}
\setbeamertemplate{navigation symbols}{}
\setbeamerfont{block title}{size=\normalsize}

\AtBeginSection[]{
  \begin{frame}[plain]
    \vfill
    \centering
    \begin{beamercolorbox}[sep=10pt,center,rounded=true,shadow=true]{palette primary}
      \usebeamerfont{title}\insertsectionhead
    \end{beamercolorbox}
    \vfill
  \end{frame}
}

% ---------------------------------------------------------------
% TikZ styles.
%   sense / plan / act  - the three stages of the autonomy loop, colour-coded
%                         once here and used identically in every deck, so a
%                         student reads the colour before reading the label.
%   node2 / topicbox    - ROS 2 computation-graph drawings.
\tikzset{
  sense/.style={rectangle, rounded corners=4pt, minimum width=2.5cm,
                minimum height=0.8cm, text centered, draw=accent!80!black,
                fill=accent!15, font=\footnotesize},
  plan/.style={rectangle, rounded corners=4pt, minimum width=2.5cm,
               minimum height=0.8cm, text centered, draw=secondary,
               fill=secondary!12, font=\footnotesize},
  act/.style={rectangle, rounded corners=4pt, minimum width=2.5cm,
              minimum height=0.8cm, text centered, draw=success!70!black,
              fill=success!18, font=\footnotesize},
  world/.style={ellipse, minimum width=2.4cm, minimum height=0.9cm,
                text centered, draw=primary, fill=primary!10, font=\footnotesize},
  rosnode/.style={ellipse, minimum width=2.4cm, minimum height=0.85cm,
                  text centered, draw=secondary, fill=secondary!12,
                  font=\footnotesize\ttfamily},
  topicbox/.style={rectangle, minimum width=2.0cm, minimum height=0.55cm,
                   text centered, draw=primary!60, fill=softgrey,
                   font=\scriptsize\ttfamily},
  hw/.style={rectangle, minimum width=2.2cm, minimum height=0.8cm,
             text centered, draw=warning!80!black, fill=warning!20,
             font=\footnotesize},
  flowarrow/.style={-{Stealth[length=2.5mm]}, thick, draw=primary},
  dataflow/.style={-{Stealth[length=2mm]}, semithick, draw=secondary}
}

% ---------------------------------------------------------------
% Full-bleed infographic slide. The title bar is part of the image,
% so the frame carries no beamer title and no furniture.
\newcommand{\fullslide}[1]{%
  \begin{frame}[plain]
    \begin{tikzpicture}[remember picture,overlay]
      \node at (current page.center)
        {\includegraphics[width=\paperwidth,height=\paperheight]{#1}};
    \end{tikzpicture}
  \end{frame}%
}

% ---------------------------------------------------------------
% RULE: every code box in this course wears the macOS terminal window bar.
% The three traffic lights are what tell a student, at a glance, that they are
% looking at something that RUNS ON A MACHINE rather than at prose about it.
% There is no bare-box variant: the title keys are applied AFTER the caller's
% options, so passing [notitle] has no effect. Do not add an exception.
%
%   \begin{pycode} ... \end{pycode}   a Python (rclpy) source file
%   \begin{shell}  ... \end{shell}    a terminal session or program output
%   \begin{spec}   ... \end{spec}     NOT runnable: .msg definitions, URDF
%                                     fragments, pseudocode, mathematics
%
% The `...s' variants (pycodes, shells) are the small ones, for a crowded slide.
%
% STANDING CONSTRAINT -- read this before writing a .msg or a shell comment.
% Inside ANY of these boxes, a `#' MUST NOT be the first non-space character of a
% line. Beamer's [fragile] scanner re-reads the frame body and treats a leading
% `#' as a macro parameter, which fails with
%       ! Illegal parameter number in definition of \next.
% Leading whitespace does not help -- it is stripped before the scan.
% A `#' anywhere else on the line is fine, which is why inline comments work:
%       self.create_timer(0.1, self.tick)   # 10 Hz, forever      <- OK
%       # geometry_msgs/msg/Twist                                 <- BREAKS THE BUILD
% For a header comment, put the file or type name in the prose ABOVE the box.
\usepackage{tcolorbox}
\tcbuselibrary{listings,skins}
\usepackage{listings}

\definecolor{termbg}{RGB}{30,30,32}
\definecolor{termbar}{RGB}{58,58,60}
\definecolor{termedge}{RGB}{88,88,92}
\definecolor{termfg}{RGB}{219,219,219}
\definecolor{termcom}{RGB}{106,153,85}
\definecolor{termkw}{RGB}{86,156,214}
\definecolor{termstr}{RGB}{206,145,120}
\definecolor{termfun}{RGB}{220,220,170}
\definecolor{termerr}{RGB}{241,76,76}
\definecolor{dotr}{RGB}{255,95,86}
\definecolor{doty}{RGB}{255,189,46}
\definecolor{dotg}{RGB}{39,201,63}

\newcommand{\termdots}{%
  \begin{tikzpicture}[baseline=-0.65ex]
    \fill[dotr] (0,0)     circle (2.7pt);
    \fill[doty] (0.34,0)  circle (2.7pt);
    \fill[dotg] (0.68,0)  circle (2.7pt);
  \end{tikzpicture}}

% Python source: VS Code "Dark+" palette, so what a student sees on the slide
% is what their editor shows them ten minutes later.
\lstdefinestyle{darkpy}{
  language=Python,
  basicstyle=\ttfamily\footnotesize\color{termfg},
  keywordstyle=\color{termkw},
  stringstyle=\color{termstr},
  commentstyle=\color{termcom},
  emph={print,len,range,float,int,str,super,min,max,abs,round},
  emphstyle=\color{termfun},
  emph={[2]rclpy,Node,create_publisher,create_subscription,create_timer,spin,
        Twist,LaserScan,Odometry,String},
  emphstyle={[2]\color{termfun}},
  showstringspaces=false, columns=fullflexible, keepspaces=true,
}
% Terminal output: same lexer, keywords left plain, errors picked out in red.
\lstdefinestyle{darkterm}{
  language=Python,
  basicstyle=\ttfamily\footnotesize\color{termfg},
  keywordstyle=\color{termfg},
  stringstyle=\color{termstr},
  commentstyle=\color{termcom},
  emph={ros2,colcon,rviz2,gazebo,source,export},
  emphstyle=\color{termfun},
  emph={[2]ERROR,error,FATAL,Failed,failed,Traceback},
  emphstyle={[2]\color{termerr}\bfseries},
  emph={[3]INFO,OK,PASS},
  emphstyle={[3]\color{termcom}},
  showstringspaces=false, columns=fullflexible, keepspaces=true,
}

\tcbset{termbox/.style={
  enhanced, listing only,
  colback=termbg, colframe=termedge, boxrule=0.5pt,
  arc=9pt, outer arc=9pt,
  left=10pt, right=10pt, top=6pt, bottom=6pt,
  title={\termdots}, colbacktitle=termbar, coltitle=termfg,
  toptitle=3.5pt, bottomtitle=3.5pt, lefttitle=10pt, titlerule=0pt,
}}
\lstdefinestyle{darkpys}{style=darkpy,     basicstyle=\ttfamily\scriptsize\color{termfg}}
\lstdefinestyle{darkterms}{style=darkterm, basicstyle=\ttfamily\scriptsize\color{termfg}}
\tcbset{tightpad/.style={top=4pt, bottom=4pt, toptitle=3pt, bottomtitle=3pt}}

\newtcblisting{pycode}[1][]{termbox, listing style=darkpy,    #1, title={\termdots}, colbacktitle=termbar, coltitle=termfg}
\newtcblisting{shell}[1][]{termbox,  listing style=darkterm,  #1, title={\termdots}, colbacktitle=termbar, coltitle=termfg}
\newtcblisting{pycodes}[1][]{termbox, listing style=darkpys,  tightpad, #1, title={\termdots}, colbacktitle=termbar, coltitle=termfg}
\newtcblisting{shells}[1][]{termbox,  listing style=darkterms, tightpad, #1, title={\termdots}, colbacktitle=termbar, coltitle=termfg}

% The light box: for things that are NOT runnable -- .msg and .srv interface
% definitions, URDF fragments, pseudocode. The contrast with the terminal box is
% the point: no window bar means nothing here runs.
%
% NOTE: the style MUST be a named \lstdefinestyle used via `listing style='.
% Passing `listing options={...}' inline instead re-parses the body and breaks on
% any `#' in the listing -- which every .msg comment starts with.
\lstdefinestyle{specstyle}{
  basicstyle=\ttfamily\footnotesize\color{primary!85!black},
  commentstyle=\color{primary!55},
  columns=fullflexible, keepspaces=true, showstringspaces=false,
}
\newtcblisting{spec}[1][]{
  enhanced, listing only, listing style=specstyle,
  colback=softgrey, colframe=primary!35, boxrule=0.6pt, arc=3pt,
  left=8pt, right=8pt, top=5pt, bottom=5pt, #1}

% ---------------------------------------------------------------
% "Under the hood" - the machine-level aside.
%
% STANDING DEVICE. Wherever a slide states a rule about how a robot or ROS 2
% behaves, the reason belongs in one of these, immediately after it. A student
% who can answer "why does it behave that way?" with something other than
% "because ROS does" is not memorising.
%
% FOUR LINES MAXIMUM. Deliberately out of scope in this course: the DDS wire
% protocol, real-time kernel scheduling, and the internals of the executor.
% Naming the exclusion here is what stops it creeping in.
\newenvironment{underhood}[1][Under the hood]{%
  \begin{tcolorbox}[enhanced, colback=primary!6, colframe=primary!45,
                    boxrule=0.6pt, arc=3pt,
                    left=8pt, right=8pt, top=5pt, bottom=5pt,
                    title={#1}, colbacktitle=primary!85!black, coltitle=white,
                    fonttitle=\footnotesize\bfseries,
                    toptitle=2pt, bottomtitle=2pt, titlerule=0pt]%
}{\end{tcolorbox}}

% "In the field" - the industrial-reality aside. Same shape, different colour:
% where \underhood explains a mechanism, this one says what it costs in a real
% deployment. Use sparingly, and never as decoration.
\newenvironment{inthefield}[1][In the field]{%
  \begin{tcolorbox}[enhanced, colback=success!7, colframe=success!50!black,
                    boxrule=0.6pt, arc=3pt,
                    left=8pt, right=8pt, top=5pt, bottom=5pt,
                    title={#1}, colbacktitle=success!45!black, coltitle=white,
                    fonttitle=\footnotesize\bfseries,
                    toptitle=2pt, bottomtitle=2pt, titlerule=0pt]%
}{\end{tcolorbox}}

% ---------------------------------------------------------------
% \imageslide{title}{figure file}{the one sentence to say about it}
%
% A PDF shows video as a still, so the deck uses images. The clip, if there is
% one, is opened from a browser -- its URL lives in the week README, never on a
% slide, because a dead link projected in front of a class is worse than no link.
%
% RULE: an image slot carries ONE sentence, and that sentence names what the
% students are looking at in the course's own vocabulary. A picture with no
% naming sentence is decoration.
%
% Generated figures are made with ../shared/gen_figure.py and each keeps its
% prompt beside it as <name>.prompt.txt, so any figure in this course can be
% regenerated or restyled consistently.
\newcommand{\imageslide}[3]{%
  \begin{frame}{#1}
    \vspace{-0.2em}
    \begin{center}
      \includegraphics[width=0.93\textwidth,height=0.63\textheight,keepaspectratio]{#2}
    \end{center}
    \vspace{0.3em}
    \centering{\small #3}
  \end{frame}%
}

% ---------------------------------------------------------------
% THE PINNED STACK. This macro and setup/README.md are the ONLY two places in
% the course where the distribution is named. Every deck, exercise and lab
% sheet writes \rosver, never the literal name. A distribution upgrade is then
% two edits, not a search across fifteen weeks of LaTeX that will miss three.
\newcommand{\rosver}{ROS\,2 Jazzy}
\newcommand{\ubuntuver}{Ubuntu 24.04 LTS}
\newcommand{\gazebover}{Gazebo Harmonic}
\newcommand{\pyver}{3.12}

%     \title...\author...\date...
%     \begin{document} ... \end{document}
%
% Build any deck with TWO pdflatex passes (three if it uses \fullslide,
% whose TikZ page anchors come from the .aux).
%
% Adapted from the SE 100 preamble: same palette, same terminal-window code
% rule, same "under the lid" device renamed \begin{underhood} for this course.
% =====================================================================

\usepackage[T1]{fontenc}
\usepackage[utf8]{inputenc}
\usepackage{graphicx}
\usepackage{amsmath}
\usepackage{amssymb}
\usepackage{tikz}
\usetikzlibrary{arrows.meta,shapes.geometric,positioning,calc,fit,backgrounds}
\usepackage{booktabs}
\usepackage{array}
\usepackage{hyperref}

\graphicspath{{./}{./figures/}}

% ---------------------------------------------------------------
% House colour palette - shared with the other Alfaisal course decks.
\definecolor{primary}{RGB}{10,37,64}
\definecolor{secondary}{RGB}{99,91,255}
\definecolor{accent}{RGB}{0,212,255}
\definecolor{success}{RGB}{50,213,131}
\definecolor{warning}{RGB}{255,193,7}
\definecolor{danger}{RGB}{220,53,69}
\definecolor{softgrey}{RGB}{240,242,245}

\setbeamercolor{structure}{fg=secondary}
\setbeamercolor{palette primary}{bg=primary,fg=white}
\setbeamercolor{palette secondary}{bg=secondary,fg=white}
\setbeamercolor{palette tertiary}{bg=accent,fg=white}
\setbeamertemplate{navigation symbols}{}
\setbeamerfont{block title}{size=\normalsize}

\AtBeginSection[]{
  \begin{frame}[plain]
    \vfill
    \centering
    \begin{beamercolorbox}[sep=10pt,center,rounded=true,shadow=true]{palette primary}
      \usebeamerfont{title}\insertsectionhead
    \end{beamercolorbox}
    \vfill
  \end{frame}
}

% ---------------------------------------------------------------
% TikZ styles.
%   sense / plan / act  - the three stages of the autonomy loop, colour-coded
%                         once here and used identically in every deck, so a
%                         student reads the colour before reading the label.
%   node2 / topicbox    - ROS 2 computation-graph drawings.
\tikzset{
  sense/.style={rectangle, rounded corners=4pt, minimum width=2.5cm,
                minimum height=0.8cm, text centered, draw=accent!80!black,
                fill=accent!15, font=\footnotesize},
  plan/.style={rectangle, rounded corners=4pt, minimum width=2.5cm,
               minimum height=0.8cm, text centered, draw=secondary,
               fill=secondary!12, font=\footnotesize},
  act/.style={rectangle, rounded corners=4pt, minimum width=2.5cm,
              minimum height=0.8cm, text centered, draw=success!70!black,
              fill=success!18, font=\footnotesize},
  world/.style={ellipse, minimum width=2.4cm, minimum height=0.9cm,
                text centered, draw=primary, fill=primary!10, font=\footnotesize},
  rosnode/.style={ellipse, minimum width=2.4cm, minimum height=0.85cm,
                  text centered, draw=secondary, fill=secondary!12,
                  font=\footnotesize\ttfamily},
  topicbox/.style={rectangle, minimum width=2.0cm, minimum height=0.55cm,
                   text centered, draw=primary!60, fill=softgrey,
                   font=\scriptsize\ttfamily},
  hw/.style={rectangle, minimum width=2.2cm, minimum height=0.8cm,
             text centered, draw=warning!80!black, fill=warning!20,
             font=\footnotesize},
  flowarrow/.style={-{Stealth[length=2.5mm]}, thick, draw=primary},
  dataflow/.style={-{Stealth[length=2mm]}, semithick, draw=secondary}
}

% ---------------------------------------------------------------
% Full-bleed infographic slide. The title bar is part of the image,
% so the frame carries no beamer title and no furniture.
\newcommand{\fullslide}[1]{%
  \begin{frame}[plain]
    \begin{tikzpicture}[remember picture,overlay]
      \node at (current page.center)
        {\includegraphics[width=\paperwidth,height=\paperheight]{#1}};
    \end{tikzpicture}
  \end{frame}%
}

% ---------------------------------------------------------------
% RULE: every code box in this course wears the macOS terminal window bar.
% The three traffic lights are what tell a student, at a glance, that they are
% looking at something that RUNS ON A MACHINE rather than at prose about it.
% There is no bare-box variant: the title keys are applied AFTER the caller's
% options, so passing [notitle] has no effect. Do not add an exception.
%
%   \begin{pycode} ... \end{pycode}   a Python (rclpy) source file
%   \begin{shell}  ... \end{shell}    a terminal session or program output
%   \begin{spec}   ... \end{spec}     NOT runnable: .msg definitions, URDF
%                                     fragments, pseudocode, mathematics
%
% The `...s' variants (pycodes, shells) are the small ones, for a crowded slide.
%
% STANDING CONSTRAINT -- read this before writing a .msg or a shell comment.
% Inside ANY of these boxes, a `#' MUST NOT be the first non-space character of a
% line. Beamer's [fragile] scanner re-reads the frame body and treats a leading
% `#' as a macro parameter, which fails with
%       ! Illegal parameter number in definition of \next.
% Leading whitespace does not help -- it is stripped before the scan.
% A `#' anywhere else on the line is fine, which is why inline comments work:
%       self.create_timer(0.1, self.tick)   # 10 Hz, forever      <- OK
%       # geometry_msgs/msg/Twist                                 <- BREAKS THE BUILD
% For a header comment, put the file or type name in the prose ABOVE the box.
\usepackage{tcolorbox}
\tcbuselibrary{listings,skins}
\usepackage{listings}

\definecolor{termbg}{RGB}{30,30,32}
\definecolor{termbar}{RGB}{58,58,60}
\definecolor{termedge}{RGB}{88,88,92}
\definecolor{termfg}{RGB}{219,219,219}
\definecolor{termcom}{RGB}{106,153,85}
\definecolor{termkw}{RGB}{86,156,214}
\definecolor{termstr}{RGB}{206,145,120}
\definecolor{termfun}{RGB}{220,220,170}
\definecolor{termerr}{RGB}{241,76,76}
\definecolor{dotr}{RGB}{255,95,86}
\definecolor{doty}{RGB}{255,189,46}
\definecolor{dotg}{RGB}{39,201,63}

\newcommand{\termdots}{%
  \begin{tikzpicture}[baseline=-0.65ex]
    \fill[dotr] (0,0)     circle (2.7pt);
    \fill[doty] (0.34,0)  circle (2.7pt);
    \fill[dotg] (0.68,0)  circle (2.7pt);
  \end{tikzpicture}}

% Python source: VS Code "Dark+" palette, so what a student sees on the slide
% is what their editor shows them ten minutes later.
\lstdefinestyle{darkpy}{
  language=Python,
  basicstyle=\ttfamily\footnotesize\color{termfg},
  keywordstyle=\color{termkw},
  stringstyle=\color{termstr},
  commentstyle=\color{termcom},
  emph={print,len,range,float,int,str,super,min,max,abs,round},
  emphstyle=\color{termfun},
  emph={[2]rclpy,Node,create_publisher,create_subscription,create_timer,spin,
        Twist,LaserScan,Odometry,String},
  emphstyle={[2]\color{termfun}},
  showstringspaces=false, columns=fullflexible, keepspaces=true,
}
% Terminal output: same lexer, keywords left plain, errors picked out in red.
\lstdefinestyle{darkterm}{
  language=Python,
  basicstyle=\ttfamily\footnotesize\color{termfg},
  keywordstyle=\color{termfg},
  stringstyle=\color{termstr},
  commentstyle=\color{termcom},
  emph={ros2,colcon,rviz2,gazebo,source,export},
  emphstyle=\color{termfun},
  emph={[2]ERROR,error,FATAL,Failed,failed,Traceback},
  emphstyle={[2]\color{termerr}\bfseries},
  emph={[3]INFO,OK,PASS},
  emphstyle={[3]\color{termcom}},
  showstringspaces=false, columns=fullflexible, keepspaces=true,
}

\tcbset{termbox/.style={
  enhanced, listing only,
  colback=termbg, colframe=termedge, boxrule=0.5pt,
  arc=9pt, outer arc=9pt,
  left=10pt, right=10pt, top=6pt, bottom=6pt,
  title={\termdots}, colbacktitle=termbar, coltitle=termfg,
  toptitle=3.5pt, bottomtitle=3.5pt, lefttitle=10pt, titlerule=0pt,
}}
\lstdefinestyle{darkpys}{style=darkpy,     basicstyle=\ttfamily\scriptsize\color{termfg}}
\lstdefinestyle{darkterms}{style=darkterm, basicstyle=\ttfamily\scriptsize\color{termfg}}
\tcbset{tightpad/.style={top=4pt, bottom=4pt, toptitle=3pt, bottomtitle=3pt}}

\newtcblisting{pycode}[1][]{termbox, listing style=darkpy,    #1, title={\termdots}, colbacktitle=termbar, coltitle=termfg}
\newtcblisting{shell}[1][]{termbox,  listing style=darkterm,  #1, title={\termdots}, colbacktitle=termbar, coltitle=termfg}
\newtcblisting{pycodes}[1][]{termbox, listing style=darkpys,  tightpad, #1, title={\termdots}, colbacktitle=termbar, coltitle=termfg}
\newtcblisting{shells}[1][]{termbox,  listing style=darkterms, tightpad, #1, title={\termdots}, colbacktitle=termbar, coltitle=termfg}

% The light box: for things that are NOT runnable -- .msg and .srv interface
% definitions, URDF fragments, pseudocode. The contrast with the terminal box is
% the point: no window bar means nothing here runs.
%
% NOTE: the style MUST be a named \lstdefinestyle used via `listing style='.
% Passing `listing options={...}' inline instead re-parses the body and breaks on
% any `#' in the listing -- which every .msg comment starts with.
\lstdefinestyle{specstyle}{
  basicstyle=\ttfamily\footnotesize\color{primary!85!black},
  commentstyle=\color{primary!55},
  columns=fullflexible, keepspaces=true, showstringspaces=false,
}
\newtcblisting{spec}[1][]{
  enhanced, listing only, listing style=specstyle,
  colback=softgrey, colframe=primary!35, boxrule=0.6pt, arc=3pt,
  left=8pt, right=8pt, top=5pt, bottom=5pt, #1}

% ---------------------------------------------------------------
% "Under the hood" - the machine-level aside.
%
% STANDING DEVICE. Wherever a slide states a rule about how a robot or ROS 2
% behaves, the reason belongs in one of these, immediately after it. A student
% who can answer "why does it behave that way?" with something other than
% "because ROS does" is not memorising.
%
% FOUR LINES MAXIMUM. Deliberately out of scope in this course: the DDS wire
% protocol, real-time kernel scheduling, and the internals of the executor.
% Naming the exclusion here is what stops it creeping in.
\newenvironment{underhood}[1][Under the hood]{%
  \begin{tcolorbox}[enhanced, colback=primary!6, colframe=primary!45,
                    boxrule=0.6pt, arc=3pt,
                    left=8pt, right=8pt, top=5pt, bottom=5pt,
                    title={#1}, colbacktitle=primary!85!black, coltitle=white,
                    fonttitle=\footnotesize\bfseries,
                    toptitle=2pt, bottomtitle=2pt, titlerule=0pt]%
}{\end{tcolorbox}}

% "In the field" - the industrial-reality aside. Same shape, different colour:
% where \underhood explains a mechanism, this one says what it costs in a real
% deployment. Use sparingly, and never as decoration.
\newenvironment{inthefield}[1][In the field]{%
  \begin{tcolorbox}[enhanced, colback=success!7, colframe=success!50!black,
                    boxrule=0.6pt, arc=3pt,
                    left=8pt, right=8pt, top=5pt, bottom=5pt,
                    title={#1}, colbacktitle=success!45!black, coltitle=white,
                    fonttitle=\footnotesize\bfseries,
                    toptitle=2pt, bottomtitle=2pt, titlerule=0pt]%
}{\end{tcolorbox}}

% ---------------------------------------------------------------
% \imageslide{title}{figure file}{the one sentence to say about it}
%
% A PDF shows video as a still, so the deck uses images. The clip, if there is
% one, is opened from a browser -- its URL lives in the week README, never on a
% slide, because a dead link projected in front of a class is worse than no link.
%
% RULE: an image slot carries ONE sentence, and that sentence names what the
% students are looking at in the course's own vocabulary. A picture with no
% naming sentence is decoration.
%
% Generated figures are made with ../shared/gen_figure.py and each keeps its
% prompt beside it as <name>.prompt.txt, so any figure in this course can be
% regenerated or restyled consistently.
\newcommand{\imageslide}[3]{%
  \begin{frame}{#1}
    \vspace{-0.2em}
    \begin{center}
      \includegraphics[width=0.93\textwidth,height=0.63\textheight,keepaspectratio]{#2}
    \end{center}
    \vspace{0.3em}
    \centering{\small #3}
  \end{frame}%
}

% ---------------------------------------------------------------
% THE PINNED STACK. This macro and setup/README.md are the ONLY two places in
% the course where the distribution is named. Every deck, exercise and lab
% sheet writes \rosver, never the literal name. A distribution upgrade is then
% two edits, not a search across fifteen weeks of LaTeX that will miss three.
\newcommand{\rosver}{ROS\,2 Jazzy}
\newcommand{\ubuntuver}{Ubuntu 24.04 LTS}
\newcommand{\gazebover}{Gazebo Harmonic}
\newcommand{\pyver}{3.12}
